% =====================================================================
%  4. Estrategia de prueba   (actividad TP5 de ISO/IEC/IEEE 29119-2)
% =====================================================================
\section{Estrategia de prueba}
\label{sec:estrategia}

\begin{notanormativa}[Correspondencia con el modelo de procesos]
Este apartado desarrolla la actividad \textbf{TP5 «Diseñar la estrategia de prueba»} de
\normap{2}: seleccionar niveles, tipos y técnicas coherentes con los riesgos; definir criterios
de entrada, salida, suspensión y reanudación; y establecer el enfoque de repetición y
regresión, los requisitos de entorno y las métricas. Cada decisión de este apartado debería
poder justificarse con un riesgo del apartado~\ref{sec:riesgos}.
\end{notanormativa}

% ---------------------------------------------------------------------
\subsection{Niveles de prueba}
\label{sec:niveles}

\begin{instruccion}
Identifique los niveles (subprocesos) de prueba que se ejecutarán, en qué consisten y en qué
orden se realizarán.

\textbf{Corrección respecto de la plantilla original:} la redacción anterior sugería que los
cuatro niveles habituales debían ejecutarse siempre. No es así: \emph{elija y justifique} los
niveles aplicables a su proyecto. Un nivel que no se ejecuta debe declararse como no aplicable
con su motivo, no omitirse en silencio. Separe con claridad componente, integración, sistema y
aceptación: no son sinónimos ni etapas intercambiables.
\end{instruccion}

\begin{table}[H]
\centering
\caption{Niveles de prueba seleccionados, con su justificación y orden de realización.}
\label{tab:niveles}
\begin{tabularx}{\textwidth}{|L{2.6cm}|C{1.9cm}|Y|C{1.5cm}|C{2.2cm}|}
  \hline
  \filaenc \textbf{Nivel} & \textbf{¿Se aplica?} & \textbf{En qué consiste en este proyecto / por qué se aplica o no} & \textbf{Orden} & \textbf{Objetos (\ref{tab:objetos})} \\ \hline
  Componente   & \campo{Sí/No} & \campo{Justificación} & \campo{1} & \campo{OP-01} \\ \hline
  Integración  & \campo{Sí/No} & \campo{Justificación} & \campo{2} & \campo{OP-01, OP-02} \\ \hline
  Sistema      & \campo{Sí/No} & \campo{Justificación} & \campo{3} & \campo{OP-01} \\ \hline
  Aceptación   & \campo{Sí/No} & \campo{Justificación} & \campo{4} & \campo{OP-01} \\ \hline
  \campo{Otro} & \campo{Sí/No} & \campo{Justificación} & \campo{} & \campo{} \\ \hline
\end{tabularx}
\end{table}

El Cuadro~\ref{tab:niveles} declara qué niveles se ejecutan y cuáles se descartan con su
motivo; el orden indicado condiciona el calendario del apartado~\ref{sec:calendario}.

% ---------------------------------------------------------------------
\subsection{Tipos de prueba}

\begin{instruccion}
Indique los tipos de prueba que se aplicarán (funcional, rendimiento, seguridad, usabilidad,
compatibilidad, entre otros) y con qué característica del alcance no funcional
(Cuadro~\ref{tab:alcance-nf}) o funcional (Cuadro~\ref{tab:alcance-f}) se corresponde cada uno.
Un tipo de prueba sin característica asociada suele indicar que el alcance está incompleto.
\end{instruccion}

\begin{table}[H]
\centering
\caption{Tipos de prueba a aplicar y su correspondencia con el alcance.}
\label{tab:tipos}
\begin{tabularx}{\textwidth}{|L{3.0cm}|C{2.4cm}|Y|C{2.3cm}|}
  \hline
  \filaenc \textbf{Tipo de prueba} & \textbf{Nivel(es)} & \textbf{Qué pretende evidenciar} & \textbf{Característica} \\ \hline
  \campo{Funcional} & \campo{Sistema} & \campo{Objetivo} & \campo{CF-01} \\ \hline
  \campo{Seguridad} & \campo{Sistema} & \campo{Objetivo} & \campo{CNF-01} \\ \hline
  & & & \\ \hline
\end{tabularx}
\end{table}

El Cuadro~\ref{tab:tipos} evita una omisión frecuente: declarar características no funcionales
en el alcance (apartado~\ref{sec:alcance}) y no prever ningún tipo de prueba capaz de
evaluarlas en alguno de los niveles seleccionados en el apartado~\ref{sec:niveles}.

% ---------------------------------------------------------------------
\subsection{Diseño e implementación: condiciones y elementos de cobertura}
\label{sec:diseno}

\begin{notanormativa}[Cadena de derivación de \normap{2}]
Las actividades \textbf{TD1–TD6} definen una cadena trazable que va de la base de prueba al
procedimiento ejecutable: identificar conjuntos de características (TD1) $\rightarrow$ derivar
condiciones de prueba (TD2) $\rightarrow$ derivar elementos de cobertura (TD3) $\rightarrow$
derivar casos de prueba (TD4) $\rightarrow$ ensamblar conjuntos (TD5) $\rightarrow$ derivar
procedimientos (TD6). La revisión señaló que la plantilla original saltaba de los elementos de
prueba directamente a los casos, perdiendo los eslabones intermedios que hacen verificable la
cobertura.
\end{notanormativa}

\begin{instruccion}
Registre aquí las condiciones de prueba derivadas de la base de prueba y de los riesgos, y los
elementos de cobertura que cada técnica genera sobre ellas. Los casos concretos se documentan
en el \ref{anx:casos}; este cuadro es el eslabón que los justifica.
\end{instruccion}

\begin{table}[H]
\centering
\caption{Condiciones de prueba y elementos de cobertura derivados.}
\label{tab:condiciones}
\begin{tabularx}{\textwidth}{|C{1.6cm}|C{1.6cm}|C{1.4cm}|Y|C{2.6cm}|}
  \hline
  \filaenc \textbf{Clave} & \textbf{Base (\ref{tab:base-prueba})} & \textbf{Riesgo} & \textbf{Condición de prueba (enunciado verificable)} & \textbf{Elementos de cobertura} \\ \hline
  \campo{COND-01} & \campo{BP-01} & \campo{RP-01} & \campo{El total se calcula correctamente para cada combinación de descuento e impuesto} & \campo{4 particiones + 6 valores límite} \\ \hline
  & & & & \\ \hline
  & & & & \\ \hline
\end{tabularx}
\end{table}

El Cuadro~\ref{tab:condiciones} materializa las actividades TD2 y TD3 y es el origen de la
columna «condición» de la matriz de trazabilidad del \ref{anx:trazabilidad}.

% ---------------------------------------------------------------------
\subsection{Técnicas de diseño de prueba}
\label{sec:tecnicas}

\begin{instruccion}
Indique y justifique las técnicas específicas que se aplicarán en cada nivel y sobre cada
objeto, según su naturaleza. \textbf{No es obligatorio usar todas las técnicas}: hay que elegir
las pertinentes y explicar su relación con los objetivos y los riesgos. Para cada técnica
seleccionada complete: objetivo, elemento al que aplica, justificación, cobertura que produce y
riesgo que ayuda a reducir.
\end{instruccion}

% =====================================================================
%  diagrams/tecnicas-prueba.tex
%  Figura: mapa de técnicas de diseño de prueba.
%
%  CORRECCIONES respecto del diagrama original («tenicas de prueba.png»):
%   1. El original fijaba «cobertura mínima: 100 %» para sentencias y
%      decisiones, atribuyéndolo implícitamente a la norma. Los
%      porcentajes se han retirado: cada técnica indica ahora el
%      ELEMENTO DE COBERTURA que produce, y el objetivo numérico se
%      declara —con su justificación— en el apartado 4.5 del plan.
%   2. Se añade la advertencia de que el mapa es un menú de selección:
%      no hay obligación de aplicar todas las técnicas.
%   3. Se añade el enfoque por palabras clave, ausente del original, con
%      su carácter opcional explícito.
% =====================================================================
\begin{figure}[htbp]
\centering
\begin{tikzpicture}[
  font=\scriptsize,
  fam/.style  ={draw=uvazul, very thick, fill=uvazul!12, rounded corners=2pt,
                text width=4.4cm, align=center, minimum height=1.0cm, inner sep=3pt},
  tec/.style  ={draw=gray!70, fill=gray!6, rounded corners=2pt,
                text width=4.4cm, align=left, minimum height=0.9cm, inner sep=4pt},
  nota/.style ={draw=polinaranja, fill=polfondo, rounded corners=2pt,
                text width=14.6cm, align=left, inner sep=5pt}
]

% ------------------------- familias (columnas) ------------------------
\node[fam] (f1) at (-5.2,0) {\textbf{Basadas en especificación}\\ \emph{caja negra}};
\node[fam] (f2) at (0,0)    {\textbf{Basadas en estructura}\\ \emph{caja blanca}};
\node[fam] (f3) at (5.2,0)  {\textbf{Basadas en experiencia}};

% ------------------------- columna 1: especificación ------------------
\node[tec] (a1) at (-5.2,-1.35) {\textbf{Particiones de equivalencia}\\ \emph{cubre:} cada clase válida e inválida};
\node[tec] (a2) at (-5.2,-2.65) {\textbf{Análisis de valores límite}\\ \emph{cubre:} bordes de cada partición};
\node[tec] (a3) at (-5.2,-3.95) {\textbf{Tablas de decisión}\\ \emph{cubre:} cada regla de la tabla};
\node[tec] (a4) at (-5.2,-5.25) {\textbf{Transiciones de estado}\\ \emph{cubre:} estados y transiciones};
\node[tec] (a5) at (-5.2,-6.55) {\textbf{Basadas en casos de uso}\\ \emph{cubre:} flujo básico y alternativos};

% ------------------------- columna 2: estructura ----------------------
\node[tec] (b1) at (0,-1.35) {\textbf{Cobertura de sentencias}\\ \emph{cubre:} cada sentencia ejecutable};
\node[tec] (b2) at (0,-2.65) {\textbf{Cobertura de decisiones}\\ \emph{cubre:} cada salida de cada decisión};
\node[tec] (b3) at (0,-3.95) {\textbf{Cobertura de condiciones}\\ \emph{cubre:} cada condición simple};
\node[tec] (b4) at (0,-5.25) {\textbf{Condiciones múltiples}\\ \emph{cubre:} combinaciones booleanas};
\node[tec] (b5) at (0,-6.55) {\textbf{Cobertura de rutas y bucles}\\ \emph{cubre:} rutas seleccionadas; 0, 1, N iteraciones};

% ------------------------- columna 3: experiencia ---------------------
\node[tec] (c1) at (5.2,-1.35) {\textbf{Pruebas exploratorias}\\ \emph{útil cuando} la base de prueba es escasa};
\node[tec] (c2) at (5.2,-2.65) {\textbf{Basadas en listas de verificación}\\ \emph{cubre:} los puntos de la lista};
\node[tec] (c3) at (5.2,-3.95) {\textbf{Pruebas por ataque}\\ \emph{útil para} riesgos de seguridad};
\node[tec] (c4) at (5.2,-5.25) {\textbf{Conjetura de errores}\\ \emph{apoyada en} defectos históricos};
\node[tec] (c5) at (5.2,-6.55) {\textbf{Palabras clave} (\normap{5})\\ \emph{enfoque opcional de implementación}};

% ------------------------------ enlaces -------------------------------
\foreach \a/\b in {f1/a1, f2/b1, f3/c1} { \draw[-{Latex[length=4pt]}, uvazul!80] (\a) -- (\b); }
\foreach \a/\b in {a1/a2, a2/a3, a3/a4, a4/a5,
                   b1/b2, b2/b3, b3/b4, b4/b5,
                   c1/c2, c2/c3, c3/c4, c4/c5} { \draw[gray!55] (\a) -- (\b); }

% ------------------------------- nota ---------------------------------
\node[nota] at (0,-8.15) {%
  \textbf{Este mapa es un menú de selección, no una lista de obligaciones.} Se eligen las
  técnicas que el riesgo justifique (Cuadro~\ref{tab:tecnicas}). Los \emph{elementos de
  cobertura} indicados son lo que cada técnica genera; el \emph{objetivo numérico} de cobertura
  no procede de la norma y debe declararse, con su justificación, en el
  apartado~\ref{sec:cobertura}. Alcanzar una cobertura determinada no demuestra la ausencia de
  defectos.};

\end{tikzpicture}
\caption[Mapa de técnicas de diseño de prueba y su elemento de cobertura]%
{Mapa de técnicas de diseño de prueba, agrupadas en tres familias, con el elemento de cobertura
que cada una produce. Versión corregida del diagrama de referencia: se han retirado los
porcentajes de cobertura mínima que se atribuían a la norma.}
\label{fig:tecnicas}
\end{figure}


La Figura~\ref{fig:tecnicas} clasifica las técnicas de diseño más habituales en tres familias
—basadas en especificación, basadas en estructura y basadas en experiencia— e indica qué tipo
de elemento de cobertura genera cada una. Sirve como menú de selección: se elige de él lo que
el riesgo justifique, no todo su contenido.

\begin{notanormativa}[Corrección respecto del diagrama original de técnicas]
El diagrama de técnicas del que parte esta plantilla presentaba «100\,\% de cobertura mínima»
para sentencias y decisiones. Ese valor \textbf{no procede de \norma}: es, a lo sumo, un
objetivo que un proyecto puede fijarse. En la Figura~\ref{fig:tecnicas} los porcentajes se han
retirado y sustituido por el \emph{elemento de cobertura} que cada técnica produce; el objetivo
numérico se declara, con su justificación, en el apartado~\ref{sec:cobertura}. Alcanzar una
cobertura determinada tampoco demuestra la ausencia de defectos.
\end{notanormativa}

\begin{table}[H]
\centering
\caption{Técnicas de diseño seleccionadas y su justificación.}
\label{tab:tecnicas}
\begin{tabularx}{\textwidth}{|L{2.9cm}|C{1.6cm}|C{1.6cm}|Y|C{1.5cm}|}
  \hline
  \filaenc \textbf{Técnica} & \textbf{Familia} & \textbf{Nivel / objeto} & \textbf{Justificación y elemento de cobertura que produce} & \textbf{Riesgo} \\ \hline
  \campo{Valores límite} & \campo{Especificación} & \campo{Componente / OP-01} & \campo{Por qué esta técnica y qué cubre} & \campo{RP-01} \\ \hline
  & & & & \\ \hline
  & & & & \\ \hline
\end{tabularx}
\notatabla{Familias: \textit{basada en especificación} (caja negra) · \textit{basada en estructura} (caja blanca) · \textit{basada en experiencia}.}
\end{table}

El Cuadro~\ref{tab:tecnicas} justifica cada técnica frente a un riesgo concreto; una técnica
sin riesgo asociado suele indicar esfuerzo no priorizado.

\begin{politicaacademica}
El curso exige someter a prueba la totalidad de los elementos con criticidad
\emph{Catastrófico} empleando pruebas basadas en casos de uso. Las demás técnicas se
seleccionan justificando su elección. Registre esa obligación en la columna de justificación
del Cuadro~\ref{tab:tecnicas} indicando que el motivo es la política del curso, para no
confundirla con una decisión basada en riesgo.
\end{politicaacademica}

\begin{instruccion}
Si emplea pruebas por palabras clave (\emph{keyword-driven testing}, \normap{5}) u otro enfoque
no listado en la Figura~\ref{fig:tecnicas}, agréguelo al Cuadro~\ref{tab:tecnicas} con la misma
justificación. Su ausencia del diagrama no lo excluye.
\end{instruccion}

% ---------------------------------------------------------------------
\subsection{Definición de la cobertura}
\label{sec:cobertura}

\begin{notanormativa}[Corrección respecto de la plantilla original]
Un porcentaje de cobertura es ininteligible si no se declara sobre qué se mide. La revisión
señaló además un error frecuente: \textbf{la proporción de casos ejecutados es avance de
ejecución, no cobertura}. Ejecutar el 90\,\% de los casos no significa haber cubierto el 90\,\%
de los requisitos, de las particiones ni de las decisiones del código.
\end{notanormativa}

\begin{instruccion}
Para cada medida de cobertura que vaya a reportar, declare las siete columnas del cuadro
siguiente. Si no puede llenar el numerador y el denominador, la medida no es utilizable como
criterio de salida.
\end{instruccion}

\begin{table}[H]
\centering
\caption{Definición operativa de cada medida de cobertura utilizada.}
\label{tab:cobertura}
\begin{tabularx}{\textwidth}{|C{1.5cm}|L{2.4cm}|Y|C{1.5cm}|C{2.0cm}|}
  \hline
  \filaenc \textbf{Clave} & \textbf{Elemento medido} & \textbf{Numerador / denominador y fórmula} & \textbf{Objetivo} & \textbf{Fuente del dato} \\ \hline
  \campo{COB-01} & \campo{Requisitos} & \campo{Requisitos con al menos un caso aprobado ÷ requisitos en el alcance} & \campo{100 \%} & \campo{Anexo G} \\ \hline
  \campo{COB-02} & \campo{Particiones de equivalencia} & \campo{Particiones ejercitadas ÷ particiones identificadas} & \campo{80 \%} & \campo{Anexo A} \\ \hline
  & & & & \\ \hline
\end{tabularx}
\notatabla{Elementos posibles: requisitos · escenarios · particiones de equivalencia · valores límite · sentencias · decisiones · ramas · estados y transiciones · reglas de una tabla de decisión.}
\end{table}

\begin{instruccion}
Añada bajo el cuadro la \textbf{interpretación} de cada medida: qué conclusión permite extraer
y cuál no. Por ejemplo: «COB-01 al 100\,\% indica que ningún requisito del alcance quedó sin
ejercitar; no indica que el software esté libre de defectos ni que los casos hayan sido
adecuados.»
\end{instruccion}

\campo{Interpretación de cada medida de cobertura}

El Cuadro~\ref{tab:cobertura} convierte los porcentajes de la política académica
(apartado~\ref{sec:criticidad}) en medidas calculables y auditables. Sus objetivos se retoman
como criterios de salida en el apartado~\ref{sec:criterios-salida}.

% ---------------------------------------------------------------------
\subsection{Criterios de entrada}
\label{sec:criterios-entrada}

\begin{notanormativa}[Apartado nuevo respecto de la plantilla original]
La plantilla original carecía de criterios de entrada. Sin ellos no puede decidirse cuándo es
legítimo empezar a ejecutar, y las ejecuciones fallidas por causas evitables (versión
equivocada, entorno no listo, datos ausentes) se confunden con defectos del producto. La
actividad \textbf{TE1} de \normap{2} incluye explícitamente «verificar el cumplimiento de los
criterios de entrada» antes de ejecutar.
\end{notanormativa}

\begin{notanormativa}[Criterio de entrada $\neq$ precondición de un caso de prueba]
Un \textbf{criterio de entrada} es una condición de \emph{proceso}: habilita el inicio de un
ciclo, nivel o campaña de ejecución y se verifica una vez, por un responsable identificado. Una
\textbf{precondición} es una condición de \emph{estado del sistema} necesaria para que un caso
concreto sea ejecutable, y se establece en cada ejecución de ese caso (véase el
\ref{anx:casos}). Confundirlas hace imposible auditar por qué se inició una campaña.
\end{notanormativa}

\begin{instruccion}
Defina los criterios que deben cumplirse antes de iniciar cada ciclo o nivel de prueba, y quién
los verifica. El registro de la verificación efectiva se lleva en el \ref{anx:entorno}.
\end{instruccion}

\begin{table}[H]
\centering
\caption{Criterios de entrada para iniciar un ciclo o nivel de prueba.}
\label{tab:criterios-entrada}
\begin{tabularx}{\textwidth}{|C{1.5cm}|Y|C{2.6cm}|C{2.3cm}|}
  \hline
  \filaenc \textbf{Clave} & \textbf{Criterio: qué debe estar listo} & \textbf{Cómo se verifica (evidencia)} & \textbf{Quién verifica} \\ \hline
  CE-01 & La versión del objeto de prueba está identificada y desplegada en el entorno. & \campo{Evidencia} & \campo{Rol} \\ \hline
  CE-02 & Los casos de prueba del ciclo están especificados y revisados. & \campo{Evidencia} & \campo{Rol} \\ \hline
  CE-03 & Los procedimientos de prueba están disponibles y ordenados. & \campo{Evidencia} & \campo{Rol} \\ \hline
  CE-04 & El entorno de prueba está preparado y verificado. & \campo{Evidencia} & \campo{Rol} \\ \hline
  CE-05 & Los datos de prueba están disponibles, cargados y en su versión correcta. & \campo{Evidencia} & \campo{Rol} \\ \hline
  CE-06 & Las personas responsables de ejecutar están asignadas y disponibles. & \campo{Evidencia} & \campo{Rol} \\ \hline
  \campo{CE-07} & \campo{Criterio adicional propio del proyecto} & \campo{Evidencia} & \campo{Rol} \\ \hline
\end{tabularx}
\notatabla{Estos seis criterios son el mínimo recomendado por la revisión; amplíelos según el proyecto. Si un ciclo se inicia sin cumplir alguno, regístrelo como desviación (\ref{sec:desviaciones}).}
\end{table}

El Cuadro~\ref{tab:criterios-entrada} fija las condiciones que habilitan el inicio de la
ejecución y nombra a quién las comprueba; su verificación efectiva, ciclo por ciclo, se
documenta en el \ref{anx:entorno}.

% ---------------------------------------------------------------------
\subsection{Criterios de finalización (criterios de salida)}
\label{sec:criterios-salida}

\begin{notanormativa}[Corrección respecto de la plantilla original]
La plantilla original resolvía la finalización por cobertura alcanzada, argumentando que en el
curso no se reparan los defectos encontrados. Eso mezcla tres cosas distintas que la revisión
pide separar: (i)~haber \emph{completado las actividades} planificadas, (ii)~los
\emph{resultados obtenidos}, y (iii)~la \emph{decisión de aceptación} del producto. Este
apartado define únicamente cuándo se dan por terminadas las actividades de prueba. La decisión
de aceptación se trata en el apartado~\ref{sec:aceptacion}.
\end{notanormativa}

\begin{instruccion}
Describa las condiciones bajo las cuales se considerarán completadas las actividades de prueba.
Combine al menos un criterio de cobertura (del Cuadro~\ref{tab:cobertura}) con criterios sobre
los resultados y sobre los incidentes abiertos. Indique qué ocurre si un criterio no se alcanza:
suspender, ampliar el esfuerzo o cerrar documentando la desviación.
\end{instruccion}

\begin{table}[H]
\centering
\caption{Criterios de finalización de las actividades de prueba.}
\label{tab:criterios-salida}
\begin{tabularx}{\textwidth}{|C{1.5cm}|C{2.4cm}|Y|C{3.0cm}|}
  \hline
  \filaenc \textbf{Clave} & \textbf{Tipo de criterio} & \textbf{Enunciado y valor objetivo} & \textbf{Si no se alcanza} \\ \hline
  \campo{CS-01} & \campo{Cobertura} & \campo{COB-01 ≥ 100 \% en objetos Catastróficos} & \campo{Ampliar / documentar desviación} \\ \hline
  \campo{CS-02} & \campo{Ejecución} & \campo{Todos los procedimientos planificados ejecutados o formalmente cancelados} & \campo{Acción} \\ \hline
  \campo{CS-03} & \campo{Incidentes} & \campo{Ningún incidente de severidad crítica sin analizar} & \campo{Acción} \\ \hline
  & & & \\ \hline
\end{tabularx}
\end{table}

El Cuadro~\ref{tab:criterios-salida} distingue criterios de cobertura, de ejecución y de
incidentes, de modo que el cierre no dependa de una sola cifra. Su evaluación efectiva se
reporta en el apartado~\ref{sec:evaluacion-salida}.

% ---------------------------------------------------------------------
\subsection{Criterios de suspensión y reanudación}

\begin{instruccion}
Especifique los criterios para suspender o reanudar las actividades de prueba, quién autoriza
cada decisión y qué actividades deberán repetirse tras una suspensión. Un caso típico de
suspensión es un entorno inestable que invalida los veredictos obtenidos.
\end{instruccion}

\begin{table}[H]
\centering
\caption{Criterios de suspensión y reanudación de las actividades de prueba.}
\label{tab:suspension}
\begin{tabularx}{\textwidth}{|C{2.3cm}|Y|C{2.4cm}|C{3.0cm}|}
  \hline
  \filaenc \textbf{Situación} & \textbf{Criterio} & \textbf{Autoriza} & \textbf{Qué se repite al reanudar} \\ \hline
  Suspensión & \campo{p.\,ej. el entorno deja de ser representativo} & \campo{Rol} & \campo{Alcance de la repetición} \\ \hline
  Reanudación & \campo{Condición que debe restablecerse} & \campo{Rol} & \campo{Alcance de la repetición} \\ \hline
\end{tabularx}
\end{table}

El Cuadro~\ref{tab:suspension} evita una discusión frecuente al reanudar: qué resultados
obtenidos antes de la suspensión siguen siendo válidos y cuáles deben volver a obtenerse.

% ---------------------------------------------------------------------
\subsection{Prueba de confirmación y prueba de regresión}
\label{sec:confirmacion-regresion}

\begin{notanormativa}[Corrección: no son lo mismo]
\begin{itemize}
  \item \textbf{Prueba de confirmación} (\emph{retesting}): se vuelve a ejecutar el caso que
        reveló un defecto, sobre la versión corregida, para comprobar que \emph{ese defecto
        concreto} quedó resuelto.
  \item \textbf{Prueba de regresión}: se ejecuta un conjunto de pruebas sobre funcionalidad que
        \emph{ya funcionaba}, para detectar efectos adversos introducidos por el cambio.
\end{itemize}
La revisión detectó que el diagrama de procesos original rotulaba la actividad «verificar
correcciones» como prueba de regresión. Es prueba de confirmación. Llamar «regresión» a
cualquier verificación de una corrección oculta que la segunda actividad —comprobar que no se
rompió nada más— puede no estarse realizando. La Figura~\ref{fig:procesos} representa ambas
por separado.
\end{notanormativa}

\begin{instruccion}
Especifique, por separado, cuándo se ejecuta cada una, qué conjunto de pruebas se selecciona y
con qué criterio. Para la regresión, indique el criterio de selección: todo el conjunto,
selección por riesgo, por área afectada por el cambio, o priorización.
\end{instruccion}

\begin{table}[H]
\centering
\caption{Enfoque de prueba de confirmación y de prueba de regresión.}
\label{tab:confirmacion-regresion}
\begin{tabularx}{\textwidth}{|L{3.0cm}|Y|Y|}
  \hline
  \filaenc  & \textbf{Prueba de confirmación} & \textbf{Prueba de regresión} \\ \hline
  Qué verifica & Que el defecto reportado quedó corregido. & Que el cambio no introdujo efectos adversos en lo que ya funcionaba. \\ \hline
  Cuándo se ejecuta & \campo{Disparador} & \campo{Disparador} \\ \hline
  Qué se selecciona & \campo{Criterio de selección} & \campo{Criterio de selección} \\ \hline
  Responsable & \campo{Rol} & \campo{Rol} \\ \hline
\end{tabularx}
\end{table}

El Cuadro~\ref{tab:confirmacion-regresion} obliga a decidir las dos actividades por separado;
cada nueva ejecución resultante se registra como una ejecución independiente en el
\ref{anx:ejecuciones}, sin sobrescribir la anterior.

\begin{politicaacademica}
El curso establece que el equipo de prueba \emph{no repara} los defectos encontrados. Esto no
elimina este apartado, pero condiciona su aplicación: declare si el equipo desarrollador
entregará correcciones durante el periodo y, en caso afirmativo, cómo se confirmarán. Si no
habrá correcciones, regístrelo como limitación y deje constancia de los defectos pendientes al
cierre (apartado~\ref{sec:incidentes-abiertos}).
\end{politicaacademica}

\campo{Declare aquí si habrá correcciones durante el periodo y cómo se tratarán}

% ---------------------------------------------------------------------
\subsection{Requisitos de entorno de prueba}
\label{sec:entorno}

\begin{instruccion}
Especifique los requisitos del entorno: hardware, software, redes, herramientas de prueba,
interfaces simuladas y fichas técnicas de los equipos necesarios. Incluya el procedimiento de
instalación y configuración.

\textbf{Ampliación respecto de la plantilla original:} además de \emph{describir} el entorno,
prevea cómo se \emph{comprueba} que está listo, quién lo mantiene, cómo se controlan sus
versiones y cómo se restaura a un estado conocido entre ejecuciones. Estas tareas corresponden
a las actividades \textbf{ES1 «Establecer el entorno»} y \textbf{ES2 «Mantener el entorno»} de
\normap{2}. El registro de cada preparación concreta se lleva en el \ref{anx:entorno}.
\end{instruccion}

\begin{table}[H]
\centering
\caption{Requisitos del entorno de prueba.}
\label{tab:req-entorno}
\begin{tabularx}{\textwidth}{|C{1.5cm}|C{2.2cm}|Y|C{2.2cm}|C{2.1cm}|}
  \hline
  \filaenc \textbf{Clave} & \textbf{Componente} & \textbf{Especificación técnica requerida} & \textbf{Versión objetivo} & \textbf{Responsable} \\ \hline
  \campo{ENT-01} & \campo{Servidor de aplicación} & \campo{Especificación} & \campo{v0.0} & \campo{Rol} \\ \hline
  \campo{ENT-02} & \campo{Base de datos} & \campo{Especificación} & \campo{v0.0} & \campo{Rol} \\ \hline
  & & & & \\ \hline
\end{tabularx}
\end{table}

\begin{table}[H]
\centering
\caption{Gestión del entorno: preparación, verificación, mantenimiento y restauración.}
\label{tab:gestion-entorno}
\begin{tabularx}{\textwidth}{|L{3.4cm}|Y|C{2.6cm}|}
  \hline
  \filaenc \textbf{Aspecto} & \textbf{Enfoque acordado} & \textbf{Responsable} \\ \hline
  Preparación (ES1) & \campo{Procedimiento de instalación y configuración} & \campo{Rol} \\ \hline
  Verificación de disponibilidad & \campo{Cómo se comprueba que el entorno está listo (criterio CE-04)} & \campo{Rol} \\ \hline
  Control de versiones del entorno & \campo{Cómo se identifica la configuración usada en cada ejecución} & \campo{Rol} \\ \hline
  Mantenimiento (ES2) & \campo{Gestión de cambios durante el periodo} & \campo{Rol} \\ \hline
  Restauración entre ejecuciones & \campo{Cuándo y cómo se devuelve el entorno a un estado conocido} & \campo{Rol} \\ \hline
  Liberación al cierre & \campo{Qué se desmonta y qué se conserva (\ref{sec:archivo})} & \campo{Rol} \\ \hline
\end{tabularx}
\end{table}

El Cuadro~\ref{tab:req-entorno} describe \emph{qué} entorno se necesita y el
Cuadro~\ref{tab:gestion-entorno}, \emph{cómo se gobierna} a lo largo del esfuerzo. Sin el
segundo, los veredictos registrados no son reproducibles.

% ---------------------------------------------------------------------
\subsection{Requisitos de datos de prueba}
\label{sec:datos}

\begin{instruccion}
Especifique los datos fuente relevantes para la prueba: qué conjuntos se necesitan, de dónde
provienen, dónde se almacenan y cómo se cargan. Indique si deben anonimizarse o sustituirse por
razones de confidencialidad o protección de datos personales, y quién es responsable de ello.

\textbf{Ampliación respecto de la plantilla original:} los datos también tienen versión. Debe
poder saberse qué conjunto de datos, en qué versión, se usó en cada ejecución; de lo contrario
un resultado no puede reproducirse ni discutirse.
\end{instruccion}

\begin{table}[H]
\centering
\caption{Conjuntos de datos de prueba requeridos.}
\label{tab:datos}
\begin{tabularx}{\textwidth}{|C{1.5cm}|L{2.9cm}|Y|C{1.8cm}|C{2.4cm}|}
  \hline
  \filaenc \textbf{Clave} & \textbf{Conjunto de datos} & \textbf{Contenido, origen y forma de obtención} & \textbf{Versión} & \textbf{Tratamiento de confidencialidad} \\ \hline
  \campo{DAT-01} & \campo{Catálogo de clientes} & \campo{Origen y volumen} & \campo{v1} & \campo{Anonimizado / sintético} \\ \hline
  & & & & \\ \hline
\end{tabularx}
\notatabla{Si algún conjunto contiene datos personales reales, indique la base legal de su uso y el procedimiento de eliminación segura al cierre (\ref{sec:archivo}).}
\end{table}

El Cuadro~\ref{tab:datos} identifica y versiona los datos; su preparación efectiva, ciclo por
ciclo, se registra junto con la del entorno en el \ref{anx:entorno}.

% ---------------------------------------------------------------------
\subsection{Gestión de configuración de los activos de prueba}
\label{sec:configuracion}

\begin{notanormativa}[Ampliación respecto de la plantilla original]
La plantilla original limitaba el control de versiones al propio documento. Para que un
resultado sea reproducible debe poder identificarse la versión de \emph{todos} los activos que
intervinieron en él. La actividad \textbf{TP9} de \normap{2} exige colocar el plan bajo gestión
de configuración, y \textbf{TE3} exige registrar versión y configuración en cada ejecución.
\end{notanormativa}

\begin{instruccion}
Indique dónde se almacena cada tipo de activo, cómo se identifica su versión y quién lo
mantiene. No es necesario un repositorio distinto por tipo: basta con que el plan diga dónde
está cada cosa y cómo se nombra.
\end{instruccion}

\begin{table}[H]
\centering
\caption{Control de configuración de los activos de prueba.}
\label{tab:configuracion}
\begin{tabularx}{\textwidth}{|L{3.2cm}|Y|C{3.0cm}|C{2.1cm}|}
  \hline
  \filaenc \textbf{Activo} & \textbf{Ubicación / repositorio} & \textbf{Esquema de identificación de versión} & \textbf{Responsable} \\ \hline
  Software bajo prueba & \campo{Ubicación} & \campo{p.\,ej. etiqueta de versión o commit} & \campo{Rol} \\ \hline
  Plan de prueba & \campo{Ubicación} & \campo{Ver. del documento} & \campo{Rol} \\ \hline
  Casos de prueba & \campo{Ubicación} & \campo{Esquema} & \campo{Rol} \\ \hline
  Procedimientos & \campo{Ubicación} & \campo{Esquema} & \campo{Rol} \\ \hline
  Guiones automatizados & \campo{Ubicación} & \campo{Esquema} & \campo{Rol} \\ \hline
  Datos de prueba & \campo{Ubicación} & \campo{Esquema} & \campo{Rol} \\ \hline
  Configuración del entorno & \campo{Ubicación} & \campo{Esquema} & \campo{Rol} \\ \hline
  Evidencias y registros & \campo{Ubicación} & \campo{Esquema} & \campo{Rol} \\ \hline
\end{tabularx}
\end{table}

El Cuadro~\ref{tab:configuracion} es la referencia que hace utilizables las columnas «versión»
y «configuración» del registro de ejecución (\ref{anx:ejecuciones}) y sostiene el archivo de
activos al cierre (apartado~\ref{sec:archivo}).

% ---------------------------------------------------------------------
\subsection{Métricas a recopilar}
\label{sec:metricas}

\begin{instruccion}
Describa las métricas que se calcularán a partir de la información recopilada durante la
prueba. Para cada una indique fórmula, fuente del dato, periodicidad y —sobre todo— cómo se
interpretará el resultado. Una métrica cuya interpretación no se ha acordado de antemano
suele terminar usándose para justificar una decisión ya tomada.

Distinga las métricas de las medidas de cobertura del apartado~\ref{sec:cobertura}: la
cobertura es un tipo particular de métrica y ya está definida allí.
\end{instruccion}

\begin{table}[H]
\centering
\caption{Métricas de prueba a recopilar, con su fórmula e interpretación.}
\label{tab:metricas}
\begin{tabularx}{\textwidth}{|C{1.5cm}|C{1.7cm}|L{3.1cm}|Y|C{2.0cm}|}
  \hline
  \filaenc \textbf{Clave} & \textbf{Tipo} & \textbf{Fórmula de cálculo} & \textbf{Interpretación acordada: qué indica y qué no} & \textbf{Fuente} \\ \hline
  \campo{MET-01} & \campo{Producto} & \campo{Fórmula} & \campo{Interpretación} & \campo{Anexo D} \\ \hline
  \campo{MET-02} & \campo{Proceso} & \campo{Fórmula} & \campo{Interpretación} & \campo{Anexo C} \\ \hline
  \campo{MET-03} & \campo{Proyecto} & \campo{Fórmula} & \campo{Interpretación} & \campo{Anexo E} \\ \hline
  & & & & \\ \hline
\end{tabularx}
\end{table}

El Cuadro~\ref{tab:metricas} define las métricas antes de recolectarlas; sus valores finales se
presentan y discuten en el apartado~\ref{sec:metricas-finales}.

\begin{politicaacademica}
Como requisito de acreditación, deben incluirse al menos \textbf{dos métricas de cada tipo}:
\begin{itemize}
  \item \textbf{De proceso}: permiten identificar mejoras en la eficiencia del proceso de
        desarrollo o de prueba.
  \item \textbf{De producto}: permiten medir la calidad del software.
  \item \textbf{De proyecto}: permiten medir la eficiencia del equipo o de las herramientas.
\end{itemize}
El número mínimo es una regla del curso. Fuera de él, la cantidad de métricas debe responder a
las preguntas que los interesados necesitan responder (Cuadro~\ref{tab:interesados}), no a una
cuota.
\end{politicaacademica}

% ---------------------------------------------------------------------
\subsection{Entregables de la prueba}
\label{sec:entregables}

\begin{instruccion}
Identifique los documentos y registros que se generarán. Indique para cada uno su formato, dónde
se mantendrá y quién es responsable. Los anexos de esta plantilla ya proporcionan formatos
listos para usar; si emplea una herramienta externa (gestor de incidencias, hoja de cálculo,
repositorio), basta con indicar dónde está y quién la mantiene, tal como admite la revisión.
\end{instruccion}

\begin{table}[H]
\centering
\caption{Entregables del esfuerzo de prueba.}
\label{tab:entregables}
\begin{tabularx}{\textwidth}{|L{4.4cm}|C{2.4cm}|Y|C{1.9cm}|}
  \hline
  \filaenc \textbf{Entregable} & \textbf{Formato} & \textbf{Ubicación / herramienta} & \textbf{Responsable} \\ \hline
  Plan de prueba (Parte I) & Este documento & \campo{Ubicación} & \campo{Rol} \\ \hline
  Reporte de finalización (Parte II) & Este documento & \campo{Ubicación} & \campo{Rol} \\ \hline
  Especificación de casos de prueba & \ref{anx:casos} & \campo{Ubicación} & \campo{Rol} \\ \hline
  Especificación de procedimientos & \ref{anx:procedimientos} & \campo{Ubicación} & \campo{Rol} \\ \hline
  Registro de ejecución & \ref{anx:ejecuciones} & \campo{Ubicación} & \campo{Rol} \\ \hline
  Registro de incidentes & \ref{anx:incidentes} & \campo{Ubicación} & \campo{Rol} \\ \hline
  Informes de estado & \ref{anx:informes} & \campo{Ubicación} & \campo{Rol} \\ \hline
  Registro de preparación de entorno y datos & \ref{anx:entorno} & \campo{Ubicación} & \campo{Rol} \\ \hline
  Matriz de trazabilidad & \ref{anx:trazabilidad} & \campo{Ubicación} & \campo{Rol} \\ \hline
\end{tabularx}
\end{table}

El Cuadro~\ref{tab:entregables} reúne los productos documentales del esfuerzo de prueba; los
seis últimos no existían en la plantilla original y son los que permiten reconstruir qué se
probó, en qué condiciones y con qué resultado.

% ---------------------------------------------------------------------
\subsection{Desviaciones respecto de la estrategia acordada}
\label{sec:desviaciones}

\begin{instruccion}
Indique cualquier aspecto de este plan que se aparte de la política o estrategia de prueba
aplicable (apartado~\ref{sec:supuestos}), con su justificación y la persona que la autoriza.
Registre aquí también las adaptaciones realizadas al reutilizar esta plantilla fuera del curso:
sustitución de umbrales académicos, cambio de la autoridad de aprobación o exclusión de algún
proceso.
\end{instruccion}

\begin{table}[H]
\centering
\caption{Desviaciones respecto de la estrategia acordada.}
\label{tab:desviaciones}
\begin{tabularx}{\textwidth}{|C{1.5cm}|L{3.0cm}|Y|C{2.4cm}|C{2.5cm}|}
  \hline
  \filaenc \textbf{Clave} & \textbf{Aspecto que se desvía} & \textbf{Justificación} & \textbf{Autoriza} & \textbf{Fecha} \\ \hline
  \campo{DES-01} & \campo{Aspecto} & \campo{Motivo} & \autoridadAprobacion & \campo{dd/mm/aaaa} \\ \hline
  & & & & \\ \hline
\end{tabularx}
\end{table}

El Cuadro~\ref{tab:desviaciones} deja constancia de lo que se hizo distinto y con qué
autorización; sin él, una adaptación legítima resulta indistinguible de un incumplimiento.
